\chapter{Criterio de tabla de claims}

\section{Tabla de claims}
\begin{definition}
Se llama tabla de claims al dispositivo curatorial mediante el cual las afirmaciones derivadas de demos seleccionadas se registran bajo una forma común, comparable y trazable dentro del Tomo III.
\end{definition}

\begin{observation}
La tabla de claims no sustituye la evidencia primaria ni la demo de origen. Su función es ordenar las afirmaciones ya suficientemente traducidas para hacerlas legibles, comparables y auditables en una capa común.
\end{observation}

\section{Condición mínima de entrada}
\begin{definition}
Un claim puede entrar en la tabla cuando concurren simultáneamente:
\begin{enumerate}
  \item procede de una demo seleccionada o de una pieza metodológicamente equivalente;
  \item declara dependencia canónica visible;
  \item declara alcance, condiciones mínimas de validez y condiciones mínimas de fallo;
  \item conserva trazabilidad suficiente hacia su evidencia de origen;
  \item puede distinguirse de otros claims sin ambigüedad inmediata.
\end{enumerate}
\end{definition}

\begin{observation}
La entrada en la tabla no implica verdad final cerrada. Implica solo que el claim ya alcanzó un nivel suficiente de formulación disciplinada.
\end{observation}

\section{Campos mínimos}
\begin{definition}
Toda entrada de la tabla debe poder declararse, al menos, mediante los siguientes campos:
\begin{enumerate}
  \item identificador;
  \item claim breve;
  \item fuente o demo de origen;
  \item dependencia canónica;
  \item alcance declarado;
  \item condición mínima de validez;
  \item condición mínima de fallo;
  \item estado editorial.
\end{enumerate}
\end{definition}

\begin{observation}
Estos campos no agotan toda ficha posible, pero bastan para una primera tabla disciplinada dentro del Tomo III.
\end{observation}

\section{Tabla y comparabilidad}
\begin{definition}
La tabla de claims debe permitir comparación disciplinada entre afirmaciones sin colapsarlas en una sola clase indiferenciada.
\end{definition}

\begin{observation}
La comparabilidad aquí no exige equivalencia fuerte entre claims. Exige solo una forma común de lectura y contraste.
\end{observation}

\section{Condición mínima de no entrada}
\begin{definition}
Un claim no debe entrar todavía en la tabla cuando su demo de origen sigue siendo opaca, su alcance no puede declararse con suficiente sobriedad o su formulación depende de extrapolaciones no disciplinadas.
\end{definition}

\begin{observation}
Esta condición protege al Tomo III frente a inflación curatorial o acumulación prematura de afirmaciones débiles.
\end{observation}

\section{Cláusula de alcance}
\begin{observation}
El criterio de tabla de claims fija una primera capa común de organización para afirmaciones derivadas de demos curadas. No constituye todavía una tabla final pública ni una validación exhaustiva de todos los claims posibles del corpus.
\end{observation}

